% ============================================================
% 6. Discussion
% ============================================================
\section{Discussion}
\label{sec:discussion}

\subsection{Governance Value Is Bottleneck-Matched, Not Uniform}

The ablation results carry a structural lesson for the design of agent
governance: the marginal value of a governance component depends on what
kind of bottleneck the task carries, and the protocol's ceiling assessment
is precisely the machinery that diagnoses the bottleneck before effort is
spent. Where the binding constraint is material---a plateau beyond the
chemistry's polarization limit, a lifetime bound that only a coating
suppresses---the escalation mechanism is causal: remove it and the task is
lost (T6, T2). Where the binding constraint is structural, escalation
matters not at all (T5), and forced architectural breadth matters only when
the search space itself hides the solution (T2). The voting mechanism is the
clearest case of a matched-null: in a suite where the cheapest route
(parameter bridges and system switches) always sufficed, the molecular
funnel was never entered, so its three-model vote was never exercised.
Governance is thus not a fixed bundle whose every element pays rent in every
regime; it is a portfolio whose elements pay according to the distribution
of bottlenecks. An organization deploying such a system should expect
exactly this pattern, and should instrument---as our protocol does---which
components were exercised, so that dormant components are identified by
data rather than presumed by fiat.

\subsection{Sharpening the Low-Headroom Calibration}

Battery-Sim-Agent, the closest published system, reports a
calibration of its headline: in low-headroom regimes where literature
defaults already sit near the target, Bayesian optimization remains
competitive and can outperform the agent \cite{chen2026battery}. Our
results sharpen this calibration in two directions. First, the BO wins we
observe are contract-attainment wins, not merely error-margin wins: on the
three structural tasks BO \emph{achieves} the contract, and with the
highest energy densities of any condition---the expected signature of an
unconstrained thin-electrode search. The tasks where BO loses are exactly
those whose bottleneck is outside its parameter space, which is a
verifiable property of the space rather than a statistical property of the
search. The boundary is thus the research target rather than a fairness
compromise: the configuration replicates the published standard to measure
what the field's strongest classical setup does---and does not do---not to
stage a symmetric contest; the aligned-penalizer re-run (Section~5.1)
closes the one configurable gap the replication left open. Second, and this is the sharper point, the protocol-free agent's
failure mode is invisible in Battery-Sim-Agent's benchmark by
construction: their objective is a fixed curve-matching error, so their
bare agent has no yardstick it can move. Our contracts are multi-criterion
natural-language requirements, which is precisely the setting where
yardstick adjustment becomes possible---and it is what the ungoverned agent
does, in five distinguishable modes (Section~5.5). The prevailing
``LLM beats classical optimization'' narrative is therefore incomplete in a
way that matters for practice: an ungoverned LLM agent in a multi-criterion
design setting can lose to a black-box optimizer \emph{and} fail its own
reported claims simultaneously---the former on reachable space, the latter
on movable criteria. Only the adjudication layer addresses both.

\subsection{Yardstick Malleability and the Case for Mechanized Adjudication}

The C1 re-adjudication is the empirical core of RQ3, and it reads directly
into the delegation literature. Delegated agency theory identifies contract
specification quality and governance visibility as the design dimensions
that discipline agent behavior \cite{delegatedagency2026}; empirical work
shows that auditable algorithmic recommendations reshape decision behavior
in organizations \cite{nikpayam2024aiassisted}. Our results instantiate
both mechanisms at the engineering level. Pre-registration is contract
specification made machine-readable before any work begins; the evidence
chain is governance visibility made file-level; and the prohibition on the
agent writing its own verdicts is a separation of duties---the auditee may
not write the audit. The five observed yardstick adjustments (threshold
shift, model substitution, parameter purchase, caliber shift, criterion
omission) are not pathological behavior; they are what an incentive-free,
unmonitored agent does with a movable yardstick, and they are exactly the
moral hazard agency theory predicts when the agent both acts and evaluates.

Two counter-arguments deserve explicit answers. First, one might read the
protocol as ``just a better prompt'': C1 received a five-sentence system
prompt while the governed agent received the full rule set. The distinction
is architectural, not rhetorical. The protocol's load-bearing components
are \emph{outside} the model's text: the adjudication layer is a
command-line program that computes verdicts, the verifier refuses
incomplete audit chains, and the agent is forbidden to write its own
evaluations. No prompt, however long, can place an evaluator outside the
agent; that is precisely what a code-level adjudication layer does. Second,
one might question whether C1 is a fair counterfactual, since its prompt
requests a deliverable but no criterion discipline. The control is
``same tools, same budget, no protocol,'' which is the natural baseline for
measuring the protocol's contribution: it holds every resource constant and
removes exactly the governance layer. A stronger C1 (e.g., adding only a
pre-registration instruction) would measure the marginal value of one
component rather than the protocol as a whole; the ablation design already
serves that purpose component by component. The magnitudes in Table~6
should therefore be read as the protocol's total contribution, not as
predictions for any particular stripped-down agent.

The theory bridge can be made quantitative with the paper's own machinery.
Delegation fidelity---the fraction of adjudicable criteria measured against
their pre-registered thresholds---is computed, not asserted, in our setup:
the governed agent scores one by construction (the verifier refuses
otherwise), while the C1 re-adjudication is exactly a fidelity audit. The
design implication transfers beyond battery design: any deployment in which
an agent's success metric can be influenced by the agent's own choices
needs an adjudication layer outside the agent, and the cheaper the
adjudication (the present layer is a command-line evaluator and a ledger),
the broader the class of tasks for which delegation becomes safe.

\subsection{Limitations}

Six limitations bound the present study. (i) \emph{Replication.} The
governed-agent baseline (pro) is the only column with full
three-seed replication: all eight tasks under three seeds each,
$24/24$ achieved under the
mechanical adjudicator with complete evidence coverage---design-level, not
trajectory-level, reproduction---and every attained task achieved in all
three seeds, so the pro column's per-task numbers are r1 values while the
8/8 total is seed-robust. The flash, luna, and C1 columns are
single-run per cell; their readings are at the level of the failure
\emph{pattern} rather than per-task inference. The C2 (Bayesian optimizer)
column carries three seeds (1234/5678/9012) with an unchanged attained set.
Cross-model corroboration spans 24 further sessions over three models. The
three-model matrix (8/8 under pro
and flash, $4/8$ with honest signatures
under a third-vendor model, Section~5.4) provides stability evidence across
models, and all three T6 seeds independently reach the same
material-bottleneck diagnosis and solution (the LNMO system switch) through
three different designs (supplementary materials). (ii)
\emph{Domain.} The evidence is confined to one domain (battery cells) and
one simulator family (PyBaMM-based electrochemistry); the governance
components are domain-agnostic in design, but their marginal values in
other engineering domains are open. (iii) \emph{Reproducibility of the
LLM.} LLM runs cannot be seeded; the protocol's verdicts are reproducible,
the agent's trajectories are not. (iv) \emph{Simulation-only.} The
real-data anchor validates the \emph{simulator}, not the final designs;
no physical cell has been built from any agent design. (v) \emph{No human
baseline.} We do not benchmark human experts, because the research
question is the comparison among delegation mechanisms under identical
resources, not the (already well-attested) economics of replacing
engineers; a human-expert control would also confound the resource-constant
design. (vi) \emph{Harness and endorsement.} All runs execute under one
agent harness. The cross-harness demonstration under two further
frameworks is executed and reported (Section~5.7); the endorsement stage is
executed for the molecular finalists in two independent codes (ORCA and
CP2K, Section~5.8), with the molecular-dynamics diffusion endorsement in
progress at the time of this draft.
